\documentclass[../../main.tex]{subfiles}
\graphicspath{{\subfix{../../image/}}}

\begin{document}

\section{正则性}\label{section:ch08-3-正则性}

本节讨论能量泛函极小元的光滑性。这通常是相当困难的课题，故我们将作一系列强简化假设。今后设泛函 $I[\cdot]$ 具有形式
\begin{align}
I[w]\triangleq\int_U L(Dw)-wf\,\mathrm{d}x,\label{align:ch08-3-eq01}
\end{align}
其中 $f\in L^2(U)$。还取 $q=2$，并假设增长条件
\begin{align}
|D_pL(p)|\leqslant C(|p|+1)\qquad (p\in\mathbb{R}^n).\label{align:ch08-3-eq02}
\end{align}
于是任一极小元 $u\in\mathcal{A}$ 是 Euler--Lagrange PDE
\begin{align}
-\sum_{i=1}^n\big(L_{p_i}(Du)\big)_{x_i}=f\quad\text{在 }U\text{ 中}\label{align:ch08-3-eq03}
\end{align}
的弱解；即
\begin{align}
\int_U\sum_{i=1}^n L_{p_i}(Du)v_{x_i}\,\mathrm{d}x=\int_U fv\,\mathrm{d}x\qquad\text{对所有 }v\in H_0^1(U).\label{align:ch08-3-eq04}
\end{align}

\subsection{二阶导数估计}\label{subsection:ch08-3-1-二阶导数估计}

现要证明：若 $u\in H^1(U)$ 是非线性 PDE \eqref{align:ch08-3-eq03} 的弱解，则事实上 $u\in H^2_{\mathrm{loc}}(U)$。但为此需加强 $L$ 上的增长条件。首先假设
\begin{align}
|D^2L(p)|\leqslant C\qquad (p\in\mathbb{R}^n).\label{align:ch08-3-eq05}
\end{align}
此外假设 $L$ 一致凸，即存在常数 $\theta>0$ 使
\begin{align}
\sum_{i,j=1}^n L_{p_ip_j}(p)\xi_i\xi_j\geqslant\theta|\xi|^2\qquad (p,\xi\in\mathbb{R}^n).\label{align:ch08-3-eq06}
\end{align}
显然这是第 6 章中线性 PDE 一致椭圆性条件的某种非线性类比。因此想法是尝试利用、或至少模仿该章中的某些计算。

\begin{theorem}[极小元的二阶导数]\label{theorem:ch08-3-1-二阶导数}
(i) 设 $u\in H^1(U)$ 是非线性偏微分方程\eqref{align:ch08-3-eq03} 的弱解，其中 $L$ 满足\eqref{align:ch08-3-eq05}、\eqref{align:ch08-3-eq06}。则
\begin{align*}
u\in H^2_{\mathrm{loc}}(U).
\end{align*}
(ii) 此外若 $u\in H_0^1(U)$ 且 $\partial U$ 为 $C^2$，则
\begin{align*}
u\in H^2(U),
\end{align*}
并有估计
\begin{align*}
\|u\|_{H^2(U)}\leqslant C\|f\|_{L^2(U)}.
\end{align*}
\begin{proof}
1. 大体遵循 \S6.3.1 中线性二阶椭圆 PDE 解的局部 $H^2$ 正则性相应定理的证明。
任取开集 $V\subset \subset U$，再选开集 $W$ 使 $V\subset \subset W\subset \subset U$。选取光滑截断函数 $\zeta$ 满足
\begin{align*}
\zeta\equiv 1\text{ 在 }V\text{ 上},\quad \zeta\equiv 0\text{ 在 }\mathbb{R}^n\setminus W\text{ 上},\quad 0\leqslant\zeta\leqslant 1.
\end{align*}
取 $|h|>0$ 小，选 $k\in\{1,\cdots,n\}$，在\eqref{align:ch08-3-eq04} 中代入
\begin{align*}
v\triangleq -D_k^{-h}(\zeta^2 D_k^h u).
\end{align*}
这里采用 \S5.8.2 的记号：
\begin{align*}
D_k^h u(x)=\frac{u(x+he_k)-u(x)}{h}\qquad (x\in W).
\end{align*}
利用恒等式 $\int_U u D_k^{-h}v\,\mathrm{d}x=-\int_U vD_k^h u\,\mathrm{d}x$，推得
\begin{align}
\sum_{i=1}^n\int_U D_k^h(L_{p_i}(Du))(\zeta^2 D_k^h u)_{x_i}\,\mathrm{d}x=-\int_U fD_k^{-h}(\zeta^2 D_k^h u)\,\mathrm{d}x.\label{align:ch08-3-eq07}
\end{align}
现
\begin{align}
D_k^h L_{p_i}(Du(x))&=\frac{L_{p_i}(Du(x+he_k))-L_{p_i}(Du(x))}{h}\notag\\
&=\frac1h\int_0^1\frac{\mathrm{d}}{\mathrm{d}s}L_{p_i}(sDu(x+he_k)+(1-s)Du(x))\,\mathrm{d}s\notag\\
&=\frac1h\int_0^1\sum_{j=1}^n L_{p_ip_j}(sDu(x+he_k)+(1-s)Du(x))\notag\\
&\qquad\cdot(u_{x_j}(x+he_k)-u_{x_j}(x))\,\mathrm{d}s\notag\\
&=\sum_{j=1}^n a^{ij,h}(x)D_k^h u_{x_j}(x),\label{align:ch08-3-eq08}
\end{align}
其中
\begin{align}
a^{ij,h}(x)\triangleq\int_0^1 L_{p_ip_j}(sDu(x+he_k)+(1-s)Du(x))\,\mathrm{d}s\qquad (i,j=1,\cdots,n).\label{align:ch08-3-eq09}
\end{align}
把\eqref{align:ch08-3-eq08} 代入\eqref{align:ch08-3-eq07} 并作简单计算，得到恒等式：
\begin{align}
A_1+A_2\triangleq&\sum_{i,j=1}^n\int_U\zeta^2 a^{ij,h}D_k^h u_{x_j}D_k^h u_{x_i}\,\mathrm{d}x\notag\\
&+\sum_{i,j=1}^n\int_U a^{ij,h}D_k^h u_{x_j}D_k^h u\,2\zeta\zeta_{x_i}\,\mathrm{d}x\notag\\
=&-\int_U fD_k^{-h}(\zeta^2 D_k^h u)\,\mathrm{d}x\triangleq B.\label{align:ch08-3-eq10}
\end{align}
一致凸性条件\eqref{align:ch08-3-eq06} 蕴涵
\begin{align}
A_1\geqslant\theta\int_U\zeta^2|D_k^h Du|^2\,\mathrm{d}x.\label{align:ch08-3-eq11}
\end{align}
此外由\eqref{align:ch08-3-eq05} 见
\begin{align}
|A_2|&\leqslant C\int_W\zeta|D_k^h Du||D_k^h u|\,\mathrm{d}x\notag\\
&\leqslant\varepsilon\int_W\zeta^2|D_k^h Du|^2\,\mathrm{d}x+\frac{C}{\varepsilon}\int_W|D_k^h u|^2\,\mathrm{d}x.\label{align:ch08-3-eq12}
\end{align}
此外，如 \S6.3.1 中线性二阶椭圆 PDE 内部 $H^2$ 估计的证明，有
\begin{align*}
|B|\leqslant\varepsilon\int_U\zeta^2|D_k^h Du|^2\,\mathrm{d}x+\frac{C}{\varepsilon}\int_U f^2+|Du|^2\,\mathrm{d}x.
\end{align*}
取 $\varepsilon=\frac\theta4$，由上述关于 $A_1,A_2,B$ 的界推得估计
\begin{align*}
\int_U\zeta^2|D_k^h Du|^2\,\mathrm{d}x\leqslant C\int_W f^2+|D_k^h u|^2\,\mathrm{d}x\leqslant C\int_U f^2+|Du|^2\,\mathrm{d}x,
\end{align*}
最后一不等式由 \S5.8.2 中差分算子定理 (i) 成立。
2. 由于 $\zeta\equiv 1$ 在 $V$ 上，得
\begin{align*}
\int_V|D_k^h Du|^2\,\mathrm{d}x\leqslant C\int_U f^2+|Du|^2\,\mathrm{d}x
\end{align*}
对 $k=1,\cdots,n$ 及所有充分小的 $|h|>0$ 成立。于是 \S5.8.2 中差分算子定理 (ii) 蕴涵 $Du\in H^1(V)$，故 $u\in H^2(V)$。这对每个 $V\subset \subset U$ 成立；故 $u\in H^2_{\mathrm{loc}}(U)$。
3. 若 $u\in H_0^1(U)$ 是\eqref{align:ch08-3-eq03} 的弱解且 $\partial U$ 为 $C^2$，则可模仿 \S6.3.2 中边界 $H^2$ 正则性定理的证明，证得 $u\in H^2(U)$，并有估计
\begin{align*}
\|u\|_{H^2(U)}\leqslant C(\|f\|_{L^2(U)}+\|u\|_{H^1(U)});
\end{align*}
细节留给读者。现由\eqref{align:ch08-3-eq06} 得不等式
\begin{align*}
(DL(p)-DL(0))\cdot p\geqslant\theta|p|^2\qquad (p\in\mathbb{R}^n).
\end{align*}
然后在\eqref{align:ch08-3-eq04} 中令 $v=u$，可用此估计导出界
\begin{align*}
\|u\|_{H^1(U)}\leqslant C\|f\|_{L^2(U)},
\end{align*}
从而完成证明。
\end{proof}
\end{theorem}

\subsection{关于更高阶正则性的注记}\label{subsection:ch08-3-2-高阶正则性}

接着要证明：若 $L$ 无穷次可微，则 $u$ 亦然。与 \S6.3 中为二阶线性椭圆 PDE 建立的正则性理论类比，把上节的 $H^2_{\mathrm{loc}}$ 估计推广以在更高 Sobolev 空间 $H^k_{\mathrm{loc}}(U)$（$k=3,4,\cdots$）中获得进一步估计，似乎是自然的。

但这一方法对非线性偏微分方程\eqref{align:ch08-3-eq03} 行不通。原因如下。对于线性方程，粗略地说，可以多次对方程求导，仍得到与原方程同一般形式的线性 PDE；例如见 \S6.3.1 中高阶内部正则性定理的证明。但若多次对非线性微分方程求导，所得越来越复杂的表达式很快就无法处理。需要深刻得多的思想，其完整展开超出本书范围。不过我们仍至少概述基本方案。

首先，取检验函数 $w\in C_c^\infty(U)$，选 $k\in\{1,\cdots,n\}$，在恒等式\eqref{align:ch08-3-eq04} 中令 $v=-w_{x_k}$（为简单起见现取 $f\equiv 0$）。由于现在已知 $u\in H^2_{\mathrm{loc}}(U)$，可分部积分得
\begin{align}
\int_U\sum_{i,j=1}^n L_{p_ip_j}(Du)u_{x_kx_j}w_{x_i}\,\mathrm{d}x=0.\label{align:ch08-3-eq13}
\end{align}
接着写
\begin{align}
\widetilde{u}\triangleq u_{x_k}\label{align:ch08-3-eq14}
\end{align}
及
\begin{align}
a^{ij}\triangleq L_{p_ip_j}(Du)\qquad (i,j=1,\cdots,n).\label{align:ch08-3-eq15}
\end{align}
再固定任一 $V\subset \subset U$。经逼近后由\eqref{align:ch08-3-eq13}--\eqref{align:ch08-3-eq15} 得
\begin{align}
\int_U\sum_{i,j=1}^n a^{ij}(x)\widetilde{u}_{x_j}w_{x_i}\,\mathrm{d}x=0\label{align:ch08-3-eq16}
\end{align}
对所有 $w\in H_0^1(V)$ 成立。这是说 $\widetilde{u}\in H^1(V)$ 是线性二阶椭圆 PDE
\begin{align}
-\sum_{i,j=1}^n(a^{ij}\widetilde{u}_{x_j})_{x_i}=0\quad\text{在 }V\text{ 中}\label{align:ch08-3-eq17}
\end{align}
的弱解。

但不能直接应用 \S6.3 的正则性理论由\eqref{align:ch08-3-eq17} 断言 $\widetilde{u}$ 光滑，原因是由\eqref{align:ch08-3-eq05}、\eqref{align:ch08-3-eq15} 只能推出
\begin{align*}
a^{ij}\in L^\infty(V)\qquad (i,j=1,\cdots,n).
\end{align*}
然而一个由 De Giorgi、Moser、Nash 独立证明的深刻定理断言：\eqref{align:ch08-3-eq17} 的任一弱解事实上对某个指数 $\gamma>0$ 局部 Hölder 连续。（见 Gilbarg--Trudinger [G-T, 第 8 章]。）于是若 $W\subset \subset V$，有 $\widetilde{u}\in C^{0,\gamma}(W)$，从而
\begin{align*}
u\in C^{1,\gamma}_{\mathrm{loc}}(U).
\end{align*}
回到定义\eqref{align:ch08-3-eq15}。若 $L$ 光滑，现已知 $a^{ij}\in C^{0,\gamma}_{\mathrm{loc}}(U)$（$i,j=1,\cdots,n$）。则\eqref{align:ch08-3-eq03} 与 Schauder 的一个较早定理 [G-T, 第 4、6 章] 断言事实上
\begin{align*}
u\in C^{2,\gamma}_{\mathrm{loc}}(U).
\end{align*}
但随后 $a^{ij}\in C^{1,\gamma}_{\mathrm{loc}}(U)$，故另一版本的 Schauder 估计蕴涵
\begin{align*}
u\in C^{3,\gamma}_{\mathrm{loc}}(U).
\end{align*}
可以继续这一所谓"bootstrap"论证，最终推出对 $k=1,2,\cdots$，$u\in C^{k,\gamma}_{\mathrm{loc}}(U)$，故 $u\in C^\infty(U)$。

关于变分法中正则性理论的更多内容，见 Giaquinta [Gi]。

\end{document}
